\documentclass[11pt]{article}
% The preamble every project's paper shares: % The preamble every project's paper shares: % The preamble every project's paper shares: \input{../../preamble} from
% projects/<project>/paper/. Boards come from chessboard and xskak, plots from
% pgfplots; every number a paper prints comes from its generated.tex.
\usepackage[LSBC4,T1]{fontenc}
\usepackage[utf8]{inputenc}
\usepackage{lmodern}
\usepackage[margin=1in]{geometry}
\usepackage{microtype}
\usepackage{amsmath}
\usepackage{booktabs}
\usepackage{tabularx}
\usepackage{array}
\usepackage{xcolor}
\usepackage[obeyspaces,spaces]{url}
\usepackage{xskak}
\usepackage{pgfplots}
\pgfplotsset{compat=1.18}
\usepackage{caption}
\usepackage{subcaption}
\usepackage{enumitem}
\usepackage[section]{placeins}
\usepackage[hidelinks]{hyperref}

\setchessboard{boardfontencoding=LSBC4, boardfontfamily=skaknew, blackfieldcolor=black!22,
  whitefieldcolor=white, setfontcolors, showmover=false, label=true, labelfontsize=6pt,
  boardfontsize=14pt}

% A move or line in UCI, and a position in FEN, both breakable at a slash.
\newcommand{\uci}[1]{\texttt{#1}}
\newcommand{\fen}[1]{{\footnotesize\path{#1}}}
\newcolumntype{L}{>{\raggedright\arraybackslash}X}

\setlength{\parskip}{0.4em}
\setlength{\parindent}{0pt}
 from
% projects/<project>/paper/. Boards come from chessboard and xskak, plots from
% pgfplots; every number a paper prints comes from its generated.tex.
\usepackage[LSBC4,T1]{fontenc}
\usepackage[utf8]{inputenc}
\usepackage{lmodern}
\usepackage[margin=1in]{geometry}
\usepackage{microtype}
\usepackage{amsmath}
\usepackage{booktabs}
\usepackage{tabularx}
\usepackage{array}
\usepackage{xcolor}
\usepackage[obeyspaces,spaces]{url}
\usepackage{xskak}
\usepackage{pgfplots}
\pgfplotsset{compat=1.18}
\usepackage{caption}
\usepackage{subcaption}
\usepackage{enumitem}
\usepackage[section]{placeins}
\usepackage[hidelinks]{hyperref}

\setchessboard{boardfontencoding=LSBC4, boardfontfamily=skaknew, blackfieldcolor=black!22,
  whitefieldcolor=white, setfontcolors, showmover=false, label=true, labelfontsize=6pt,
  boardfontsize=14pt}

% A move or line in UCI, and a position in FEN, both breakable at a slash.
\newcommand{\uci}[1]{\texttt{#1}}
\newcommand{\fen}[1]{{\footnotesize\path{#1}}}
\newcolumntype{L}{>{\raggedright\arraybackslash}X}

\setlength{\parskip}{0.4em}
\setlength{\parindent}{0pt}
 from
% projects/<project>/paper/. Boards come from chessboard and xskak, plots from
% pgfplots; every number a paper prints comes from its generated.tex.
\usepackage[LSBC4,T1]{fontenc}
\usepackage[utf8]{inputenc}
\usepackage{lmodern}
\usepackage[margin=1in]{geometry}
\usepackage{microtype}
\usepackage{amsmath}
\usepackage{booktabs}
\usepackage{tabularx}
\usepackage{array}
\usepackage{xcolor}
\usepackage[obeyspaces,spaces]{url}
\usepackage{xskak}
\usepackage{pgfplots}
\pgfplotsset{compat=1.18}
\usepackage{caption}
\usepackage{subcaption}
\usepackage{enumitem}
\usepackage[section]{placeins}
\usepackage[hidelinks]{hyperref}

\setchessboard{boardfontencoding=LSBC4, boardfontfamily=skaknew, blackfieldcolor=black!22,
  whitefieldcolor=white, setfontcolors, showmover=false, label=true, labelfontsize=6pt,
  boardfontsize=14pt}

% A move or line in UCI, and a position in FEN, both breakable at a slash.
\newcommand{\uci}[1]{\texttt{#1}}
\newcommand{\fen}[1]{{\footnotesize\path{#1}}}
\newcolumntype{L}{>{\raggedright\arraybackslash}X}

\setlength{\parskip}{0.4em}
\setlength{\parindent}{0pt}

% Written by projects/central-strike/paper/generate.mjs; do not edit by hand.

\newcommand{\AppCommit}{f4bd63f19b}
\newcommand{\GradeNodes}{250,000}
\newcommand{\Window}{3}
\newcommand{\Draws}{2,000}
\newcommand{\Seed}{20261020}
\newcommand{\MinGames}{30}
\newcommand{\MonthName}{September 2026}
\newcommand{\Floor}{2200}
\newcommand{\WinLow}{10}
\newcommand{\WinHigh}{90}
\newcommand{\CentreFirst}{116}
\newcommand{\AwayFirst}{225}
\newcommand{\CentreGamesFirst}{47}
\newcommand{\AwayGamesFirst}{67}
\newcommand{\CentreSecond}{83}
\newcommand{\AwaySecond}{191}
\newcommand{\CentreGamesSecond}{39}
\newcommand{\AwayGamesSecond}{59}
\newcommand{\StrongCentre}{129}
\newcommand{\StrongAway}{428}
\newcommand{\HisScored}{909}
\newcommand{\StrongScored}{1,200}
\newcommand{\HisEndgames}{90}
\newcommand{\HisNoCentral}{204}
\newcommand{\StrongEndgames}{195}
\newcommand{\StrongNoCentral}{448}
\newcommand{\Entering}{615}
\newcommand{\StrongEntering}{557}
\newcommand{\CoverageTable}{%
\begin{table}[ht]
\centering\footnotesize
\begin{tabular}{@{}lrrr@{}}
\toprule
 & \begin{tabular}[b]{@{}l@{}}His, first \\ half\end{tabular} & \begin{tabular}[b]{@{}l@{}}His, second \\ half\end{tabular} & \begin{tabular}[b]{@{}l@{}}Strong \\ players\end{tabular} \\
\midrule
King in the centre: positions (games) & 116 (47) & 83 (39) & 129 \\
King away: positions (games) & 225 (67) & 191 (59) & 428 \\
King on its starting square & 84 & 75 & 114 \\
\bottomrule
\end{tabular}
\caption{The entering positions: quiet middlegame positions with a central break on, every score the study reads exact. His two halves were counted before the rule was fixed, with no score read; of all the scored positions, his 90 endgames and 204 positions with no central break on stay out, and the strong players' 195 and 448.}
\label{tab:coverage}
\end{table}
}
\newcommand{\LiftFirst}{13.0}
\newcommand{\LiftFirstCi}{0.0 to 25.2}
\newcommand{\GainFirst}{28.4}
\newcommand{\GainFirstCi}{$-$38.7 to $-$18.1}
\newcommand{\RefereeFirst}{37.2}
\newcommand{\NullFirst}{$-$3.5}
\newcommand{\NullFirstCi}{$-$13.7 to 6.4}
\newcommand{\LiftSecond}{23.0}
\newcommand{\LiftSecondCi}{4.7 to 42.4}
\newcommand{\GainSecond}{26.5}
\newcommand{\GainSecondCi}{$-$43.5 to $-$9.2}
\newcommand{\RefereeSecond}{38.7}
\newcommand{\NullSecond}{$-$0.8}
\newcommand{\NullSecondCi}{$-$11.4 to 10.7}
\newcommand{\StrongLift}{5.8}
\newcommand{\StrongLiftCi}{$-$3.5 to 15.5}
\newcommand{\StrongGain}{58.1}
\newcommand{\StrongGainCi}{$-$67.9 to $-$48.3}
\newcommand{\StrongReferee}{54.8}
\newcommand{\StrongReading}{silent}
\newcommand{\Verdict}{no signal}
\newcommand{\RuleTable}{%
\begin{table}[ht]
\centering\footnotesize
\begin{tabular}{@{}llll@{}}
\toprule
 & \begin{tabular}[b]{@{}l@{}}His, first \\ half\end{tabular} & \begin{tabular}[b]{@{}l@{}}His, second \\ half\end{tabular} & \begin{tabular}[b]{@{}l@{}}Strong \\ players\end{tabular} \\
\midrule
Referee: $A - R$ ($S - R$) & 37.2 {\scriptsize(30.6, 43.3)} & 38.7 {\scriptsize(31.2, 45.7)} & 54.8 {\scriptsize(50.4, 59.2)} \\
Null: $R_2 - R$, centre minus away & $-$3.5 {\scriptsize($-$13.7, 6.4)} & $-$0.8 {\scriptsize($-$11.4, 10.7)} & 5.0 {\scriptsize($-$3.7, 13.8)} \\
The lift: $C - R$, centre minus away & 13.0 {\scriptsize(0.0, 25.2)} & 23.0 {\scriptsize(4.7, 42.4)} & 5.8 {\scriptsize($-$3.5, 15.5)} \\
The strike's gain: $C - A$ ($S$), centre & $-$28.4 {\scriptsize($-$38.7, $-$18.1)} & $-$26.5 {\scriptsize($-$43.5, $-$9.2)} & $-$58.1 {\scriptsize($-$67.9, $-$48.3)} \\
\bottomrule
\end{tabular}
\caption{The rule's measures, in points of soundness with 95\% intervals from 2,000 bootstrap draws over games. The verdict rests on his halves; the strong players' column is read, not ruled on.}
\label{tab:rule}
\end{table}
}
\newcommand{\ArmTable}{%
\begin{table}[ht]
\centering\footnotesize
\setlength{\tabcolsep}{3pt}
\begin{tabular}{@{}llllll@{}}
\toprule
 & \begin{tabular}[b]{@{}l@{}}First half, \\ centre\end{tabular} & \begin{tabular}[b]{@{}l@{}}First half, \\ away\end{tabular} & \begin{tabular}[b]{@{}l@{}}Second half, \\ centre\end{tabular} & \begin{tabular}[b]{@{}l@{}}Second half, \\ away\end{tabular} & \begin{tabular}[b]{@{}l@{}}Strong, \\ centre / away\end{tabular} \\
\midrule
$C$, the central breaks & 31.0 {\scriptsize(21.2, 41.3)} & 13.8 {\scriptsize(9.2, 19.2)} & 30.1 {\scriptsize(19.8, 41.8)} & 15.2 {\scriptsize(8.8, 23.0)} & 19.4 / 18.2 \\
$A$, his move ($S$) & 59.5 {\scriptsize(49.5, 69.2)} & 52.9 {\scriptsize(45.7, 60.2)} & 56.6 {\scriptsize(45.9, 68.9)} & 59.7 {\scriptsize(51.2, 67.8)} & 77.5 / 68.7 \\
$N$, Maia-2 at his band & 69.0 {\scriptsize(59.4, 78.5)} & 66.2 {\scriptsize(59.8, 73.0)} & 72.3 {\scriptsize(62.6, 82.1)} & 72.8 {\scriptsize(65.0, 79.9)} & -- \\
$U$, Maia-2 at the band above & 69.0 {\scriptsize(59.5, 78.6)} & 67.6 {\scriptsize(60.6, 74.2)} & 71.1 {\scriptsize(61.2, 80.7)} & 74.3 {\scriptsize(66.7, 81.3)} & -- \\
$R$, a random move & 20.7 {\scriptsize(13.1, 29.5)} & 16.4 {\scriptsize(11.3, 22.0)} & 14.5 {\scriptsize(6.7, 23.2)} & 22.5 {\scriptsize(15.1, 30.1)} & 12.4 / 17.1 \\
$R_2$, another & 19.8 {\scriptsize(12.4, 29.3)} & 19.1 {\scriptsize(13.8, 25.2)} & 12.0 {\scriptsize(5.2, 20.0)} & 20.9 {\scriptsize(13.8, 27.9)} & 16.3 / 15.9 \\
He played a central break & 8.6 {\scriptsize(3.6, 14.3)} & 7.1 {\scriptsize(4.0, 11.1)} & 7.2 {\scriptsize(2.6, 12.0)} & 4.7 {\scriptsize(2.0, 7.5)} & 8.5 / 11.0 \\
\bottomrule
\end{tabular}
\caption{Described, not tested: the share of each arm that is sound, and how often the move played was a central break, in percent with 95\% intervals, where the other king is in the centre and where it is away.}
\label{tab:arms}
\end{table}
}
\newcommand{\CCentreFirst}{31}
\newcommand{\CAwayFirst}{14}
\newcommand{\ACentreFirst}{59}
\newcommand{\AAwayFirst}{53}
\newcommand{\RCentreFirst}{21}
\newcommand{\RAwayFirst}{16}
\newcommand{\PlayedCentreFirst}{9}
\newcommand{\PlayedAwayFirst}{7}
\newcommand{\LostCCentreFirst}{13.4}
\newcommand{\LostCAwayFirst}{18.6}
\newcommand{\LostACentreFirst}{4.4}
\newcommand{\StrictFirst}{11.0}
\newcommand{\StrictFirstCi}{$-$5.4 to 28.2}
\newcommand{\QueensFirst}{14.5}
\newcommand{\QueensFirstCi}{0.8 to 27.1}
\newcommand{\FlankFirst}{6.7}
\newcommand{\FlankFirstCi}{$-$4.2 to 17.8}
\newcommand{\MaiaCFirst}{4.9}
\newcommand{\MaiaFFirst}{3.8}
\newcommand{\MaiaBothFirst}{206}
\newcommand{\CCentreSecond}{30}
\newcommand{\CAwaySecond}{15}
\newcommand{\ACentreSecond}{57}
\newcommand{\AAwaySecond}{60}
\newcommand{\RCentreSecond}{14}
\newcommand{\RAwaySecond}{23}
\newcommand{\PlayedCentreSecond}{7}
\newcommand{\PlayedAwaySecond}{5}
\newcommand{\LostCCentreSecond}{13.5}
\newcommand{\LostCAwaySecond}{17.9}
\newcommand{\LostACentreSecond}{5.7}
\newcommand{\StrictSecond}{19.3}
\newcommand{\StrictSecondCi}{0.5 to 38.0}
\newcommand{\QueensSecond}{23.3}
\newcommand{\QueensSecondCi}{2.8 to 42.3}
\newcommand{\FlankSecond}{5.0}
\newcommand{\FlankSecondCi}{$-$8.3 to 18.7}
\newcommand{\MaiaCSecond}{4.8}
\newcommand{\MaiaFSecond}{4.3}
\newcommand{\MaiaBothSecond}{128}
\newcommand{\StrongStrict}{6.7}
\newcommand{\StrongStrictCi}{$-$3.4 to 16.9}
\newcommand{\StrongQueens}{4.8}
\newcommand{\StrongQueensCi}{$-$6.8 to 15.4}
\newcommand{\StrongFlank}{18.0}
\newcommand{\StrongFlankCi}{7.5 to 28.1}
\newcommand{\StrongCCentre}{19}
\newcommand{\StrongCAway}{18}
\newcommand{\StrongSCentre}{78}
\newcommand{\EaFen}{\fen{r3k2r/pbp2ppp/1pnqpn2/3p4/3P4/2PBPNP1/PP1N1PP1/R2QK2R w KQkq - 2 10}}
\newcommand{\EaMover}{White}
\newcommand{\EaStructure}{the London pyramid}
\newcommand{\EaKing}{e8}
\newcommand{\EaBreak}{10.~e4}
\newcommand{\EaBreakSound}{sound}
\newcommand{\EaWpBreak}{55}
\newcommand{\EaBest}{10.~\symqueen{}a4}
\newcommand{\EaWpBest}{55}
\newcommand{\EaPlayed}{10.~\symknight{}e5}
\newcommand{\EaWpPlayed}{38}
\newcommand{\EaPlayedSound}{not sound}
\newcommand{\EaBoard}{%
\chessboard[setfen={r3k2r/pbp2ppp/1pnqpn2/3p4/3P4/2PBPNP1/PP1N1PP1/R2QK2R w KQkq - 2 10}, boardfontsize=14pt, pgfstyle=straightmove, arrow=to, linewidth=0.07em, color=blue!70!black, markmoves={f3-e5}, pgfstyle=straightmove, arrow=to, linewidth=0.07em, color=green!50!black, markmoves={e3-e4}]
}
\newcommand{\EbFen}{\fen{r2qk2r/pp3p2/3bpp1p/3p3Q/3P4/2P5/PP3PPP/RN2K2R b KQkq - 1 13}}
\newcommand{\EbMover}{Black}
\newcommand{\EbStructure}{the reversed Carlsbad}
\newcommand{\EbKing}{e1}
\newcommand{\EbBreak}{13\ldots{}e5}
\newcommand{\EbBreakSound}{not sound}
\newcommand{\EbWpBreak}{45}
\newcommand{\EbBest}{13\ldots{}\symqueen{}b6}
\newcommand{\EbWpBest}{55}
\newcommand{\EbPlayed}{13\ldots{}\symqueen{}b6}
\newcommand{\EbWpPlayed}{56}
\newcommand{\EbPlayedSound}{sound}
\newcommand{\EbBoard}{%
\chessboard[setfen={r2qk2r/pp3p2/3bpp1p/3p3Q/3P4/2P5/PP3PPP/RN2K2R b KQkq - 1 13}, boardfontsize=14pt, pgfstyle=straightmove, arrow=to, linewidth=0.07em, color=blue!70!black, markmoves={d8-b6}, pgfstyle=straightmove, arrow=to, linewidth=0.07em, color=red!75!black, markmoves={e6-e5}]
}


\title{Is a central break right while the other king is still in the centre?\\[0.3em]
\large A coach's moment for a pawn break, on \Entering{} of his quiet middlegame positions and
\StrongEntering{} of strong players', refereed by Stockfish}
\author{Chess Coach research lab}
\date{}

\begin{document}
\maketitle

\begin{abstract}
A coach reviewing an adult beginner's rapid games teaches a moment for a pawn break: strike in the
centre while the other king has not castled. The app leaves the moment to the engine, because a
break that can be played is sound about one time in five in his quiet positions, no more often
than a random move. In the positions the structure study scored, his quiet middlegame positions
with a central break on, a central break is sound \CCentreFirst\% and \CCentreSecond\% of the time
while the other king is still in the centre, against \CAwayFirst\% and \CAwaySecond\% once it has
left. Measured against a random move in the same positions, the lift is \LiftFirst{} points
(\LiftFirstCi) on the earlier half of his games and \LiftSecond{} (\LiftSecondCi) on the later
one. The earlier interval reaches 0, so by the rule fixed before any score was read by the king
there is no signal, though both halves point the same way. STRONGABSTRACT Even with the king in
the centre, the central break trails his own move by \GainFirst{} and \GainSecond{} points. The
coach's moment may hold at his level, but it is not shown here, and it would not make the break a
better choice than his own move: the verdict on a break stays the engine's.
\end{abstract}

\section{The question}

A titled player reviewing the rapid games of an adult beginner~\cite{video} teaches a moment for a pawn break: strike in the centre while the other king has not castled,
because the lines the break opens lead to the king. He adds that players at that level miss the
central break because the pawn feels like part of the scenery.

The app teaches each pawn structure's breaks, and two earlier studies shaped what it says about
them. A break that can be played is sound only about one time in five in his quiet positions, no
more often than a random move, while his own move is sound more than half the time; and the
pieces being ready for the break says nothing about whether it is sound, while the app's blunder
check does. So the app's lesson tells him to check what a break leaves hanging and leaves the
moment to the engine.

The coach's moment is a reading of the board, so it can be tested the same way. In his quiet
middlegame positions with a central break on, is the break sound more often when the other king
is still in the centre than when it has left, beyond what the king's position does for any move,
and is it then as sound as his own move? If the king marks the moment, the app can name a king
still in the centre as the time to look at the central break; if the break is then as sound as
his own move, as the time to play it. If not, the moment stays the engine's.

\section{How it was measured}

\paragraph{The positions.} The positions the structure study scored, with their scores as they
stand. His are his quiet middlegame and endgame positions with a break of his structure on, from
his 10-minute games in the app's archive at commit \texttt{\AppCommit}, in two halves by the
games' end time; nothing is fitted, so the halves are two replications. The strong players' are
positions from the Lichess games of \MonthName{} between two players rated \Floor{} or more, one a
game, the first \StrongScored{} scored in source order. A position is quiet when Stockfish's best
move is not a capture, a check or a promotion and the mover's win\% is between \WinLow{} and \WinHigh. The
referee is Stockfish at the app's grading budget of \GradeNodes{} nodes, the hash cleared before
each position; a move is sound when it gives up no more than \Window{} win\% against the best.

A position enters this study when it is a middlegame by the app's phase reading, when one of its
breaks lands on the d- or e-file (a central break), and when each central break, the move played
and the two random moves have exact scores. The other king is \emph{in the centre} when it stands
on the d-, e- or f-file, where a king that has not castled stands, and \emph{away} otherwise.

\paragraph{The measures.} In each position, $C$ is the share of its central breaks that are sound;
$A$ is 1 when his move is sound (for the strong players $S$, the strong player's move), and $R$
and $R_2$ the same for two random legal moves that are not breaks, drawn by the structure study.
The \emph{lift} is the mean of $C - R$ where the king is in the centre minus its mean where the
king is away: $R$ stands for how forgiving a position is, so the lift asks whether the king in the
centre helps the central break more than it helps any move. The \emph{strike's gain} is the mean
of $C - A$ where the king is in the centre. Intervals come from \Draws{} bootstrap draws with seed
\Seed, resampling his games within a half (a game's positions of both kinds together), or the
strong players' positions: the 95\% percentile interval.

\section{Counted before the rule}

Before the rule was fixed, his positions were counted by the king's file and the breaks' files,
with no score read (Table~\ref{tab:coverage}). The strong players' positions were being rescored
in source order at the time, after a fix to the structure study, and were counted only when the
analysis ran. His positions with the king in the centre come from \CentreGamesFirst{} and
\CentreGamesSecond{} games, so only a large lift can show on each half.

\CoverageTable

\section{The rule, fixed before any score was read by the king}

The controls come first, and if either fails there is no verdict. (1) His move must be sound more
often than $R$, by an interval above 0, on each half, and the strong player's move on theirs.
(2) The lift with $R_2$ in place of $C$ must not lie above 0 on both halves, nor below 0 on both.

Then, on his positions: the strike holds when the lift lies above 0 on both halves and the
strike's gain lies below 0 on neither; the king marks the moment but the break trails his move
when the lift lies above 0 on both halves and the gain lies below 0 on a half; the king misleads
when the lift lies below 0 on both halves; otherwise there is no signal. There is not enough data
when, on a half, the positions with the king in the centre or away come from fewer than
\MinGames{} games. The strong players' positions are read by the same measures and do not change
the verdict.

\section{The verdict: no signal}

The controls pass (Table~\ref{tab:rule}). His move is sound \RefereeFirst{} and \RefereeSecond{}
points more often than a random move, and the strong player's \StrongReferee; the null's lift is
\NullFirst{} (\NullFirstCi) and \NullSecond{} (\NullSecondCi), holding 0 on both halves.

The lift is \LiftFirst{} points (\LiftFirstCi) on the earlier half and \LiftSecond{}
(\LiftSecondCi) on the later one. Both are positive, but the earlier interval reaches 0, so by the
rule there is no signal. STRONGVERDICT

The strike's gain answers the second question whatever the first: with the king in the centre, the
central break is sound \GainFirst{} and \GainSecond{} points less often than his own move
(\GainFirstCi{} and \GainSecondCi), and \StrongGain{} less often than the strong player's.

\RuleTable

\section{What the positions describe}

None of what follows was part of the rule (Table~\ref{tab:arms}).

\paragraph{The arms.} With the king in the centre a central break is sound \CCentreFirst\% and
\CCentreSecond\% of the time, and \CAwayFirst\% and \CAwaySecond\% with it away, while a random
move is sound \RCentreFirst\% and \RCentreSecond\% in the first case and \RAwayFirst\% and
\RAwaySecond\% in the second. His own move is sound \ACentreFirst\% and \ACentreSecond\% with the
king in the centre. The central breaks give up \LostCCentreFirst{} and \LostCCentreSecond{} win\%
on average with the king in the centre, against \LostCAwayFirst{} and \LostCAwaySecond{} with it
away; his move gives up \LostACentreFirst{} and \LostACentreSecond.

\paragraph{Other readings.} With the king in the centre only when it stands on its starting square,
the lift is \StrictFirst{} (\StrictFirstCi) and \StrictSecond{} (\StrictSecondCi). With both queens
on the board, it is \QueensFirst{} (\QueensFirstCi) and \QueensSecond{} (\QueensSecondCi). For the
breaks that land off the d- and e-files, the same lift is \FlankFirst{} (\FlankFirstCi) and
\FlankSecond{} (\FlankSecondCi), holding 0 on both halves. STRONGOTHER

\paragraph{Finding it.} He played a central break in \PlayedCentreFirst\% and
\PlayedCentreSecond\% of his positions with the king in the centre, and \PlayedAwayFirst\% and
\PlayedAwaySecond\% with it away. In the positions with both kinds of break on, Maia-2 at his band
gives a central break \MaiaCFirst\% and \MaiaCSecond\% on average and a flank break \MaiaFFirst\%
and \MaiaFSecond\%: by that reading a central break is no harder for his level to see than a flank
break.

\ArmTable

\section{Two positions}

Two positions of his later half with the other king in the centre (Figure~\ref{fig:examples}): the
first where a central break is sound and he played something else, and the first where every
central break is unsound and his own move is sound. They illustrate; they test nothing.

\begin{figure}[ht]
\centering
\begin{subfigure}[t]{0.48\textwidth}
\centering
\EaBoard\\[0.3em]
{\footnotesize FEN \EaFen}
\caption{He is \EaMover, in \EaStructure; the other king is on \EaKing. Stockfish: \EaBest{} is
best at \EaWpBest\% to win; the central break \EaBreak{} is \EaBreakSound{} (\EaWpBreak\%). He
played \EaPlayed{} (blue), \EaPlayedSound{} at \EaWpPlayed\%.}
\label{fig:struck}
\end{subfigure}\hfill
\begin{subfigure}[t]{0.48\textwidth}
\centering
\EbBoard\\[0.3em]
{\footnotesize FEN \EbFen}
\caption{He is \EbMover, in \EbStructure; the other king is on \EbKing. Stockfish: \EbBest{} is
best at \EbWpBest\% to win; the central break \EbBreak{} is \EbBreakSound{} (\EbWpBreak\%). He
played \EbPlayed{} (blue), \EbPlayedSound{} at \EbWpPlayed\%.}
\label{fig:trailed}
\end{subfigure}
\caption{Two of his positions with the other king in the centre. The break is green where
Stockfish finds it sound and red where it does not; win\% is his after each move. Stockfish at
\GradeNodes{} nodes, rerun for this figure.}
\label{fig:examples}
\end{figure}

\section{What it means for the app}

The app's pawn-break lesson stays as it is: check what a break leaves hanging before playing it,
and leave the moment to the engine. A king still in the centre may be a reason to look at the
central break, and both halves of his games lean that way, but at this size the lean is not shown,
and where it shows most the break is still sound less often than his own move. A test on Lichess
games at his band, with many more positions, could settle the lean; it would not change the second
answer.

\section{Caveats}

\begin{itemize}[nosep]
\item \textbf{Few games.} His positions with the king in the centre come from a few dozen games a
  half, so only a large lift can show on each half.
\item \textbf{Not a cause.} A king in the centre comes with an earlier middlegame and more pieces;
  $R$ stands for how forgiving a position is, not for every difference.
\item \textbf{A break stands for a plan}, tested as the next move.
\item \textbf{Mostly the London.} Most of his positions are the London pyramid, whose central
  breaks are e3-e4 for White and \ldots e6-e5 for Black.
\item \textbf{Mostly blitz.} Most of the strong players' games are blitz; his are rapid.
\end{itemize}

\begin{thebibliography}{9}
\bibitem{video} A titled player's review of an adult beginner's rapid games, ``Anyone Can Be 1200
  In Chess. Watch This.'' (video, \url{https://www.youtube.com/watch?v=PcGDehRSqTM}). The lab read
  a summary of its captions, not the video itself.
\end{thebibliography}

\end{document}
